\documentclass[12pt,a4paper]{article}

\usepackage[margin=1in]{geometry}
\usepackage[hidelinks]{hyperref}
\usepackage{enumitem}
\usepackage{titlesec}

\setlength{\parindent}{0pt}
\setlength{\parskip}{8pt}

\titleformat{\section}{\large\bfseries}{}{0pt}{}
\titleformat{\subsection}{\normalsize\bfseries}{}{0pt}{}

\begin{document}

\begin{center}
    {\LARGE Portfolio Project Report}\\[6pt]
    {\large Working Student IT System Application Portfolio}\\[12pt]
    \textbf{Name:} Adwaid Sreeletha Anil\\
    \textbf{Course:} BSc Computer Science\\
    \textbf{University:} GISMA University of Applied Sciences\\
    \textbf{Email:} \href{mailto:adwaithh9383@gmail.com}{adwaithh9383@gmail.com}\\
    \textbf{Phone:} \href{tel:+4917623862783}{+49 17623862783}\\
    \textbf{Location:} Berlin, Germany
\end{center}

\section{1. Purpose of the Portfolio}

The purpose of this project is to create a personal computer science portfolio to apply for work student jobs. It shows my technical skills, CV, GitHub profile and my basic python project. Since I am still at the beginning of my computer science journey, Since i'm a complete beginner and new to computer journey, my portfolio is simple and easy to understand. This portfolio is made for recruiters, professors who can easily learn what i am and my skills so far.

\section{2.Target Student Job Link}

\textbf{Position:} Working Student- Web Development\\
\textbf{Company:} CareHelper24 UG\\
\textbf{Location:} Kantstaße 72, 10627 Berlin, Germany\\
\textbf{Job Link:} \url{https://de.indeed.com/viewjob?jk=84413f48e7e55170&from=shareddesktop_copy}

\section{3. Target Role Summary}

The target role is Working Student (m/f/d) – Web Development. The position involves developing and maintaining web applications, building frontend features using HTML, CSS, and JavaScript, improving backend functions, fixing technical issues, and supporting the development team. It is a good opportunity to gain practical experience while improving programming and web development skills.

\section{4. My Experience and Knowledge}

I am a BSc Computer Science student at GISMA University of Applied Sciences. I'm a complete beginner but i am learning from basic foundation and doing small projects and practicing daily, also try to learn new things everyday. I have basic knowledge of HTML, CSS, Python, Git, GitHub, LaTeX, JavaScript and also i"m improving my touch typing skills.

My biggest strength is my passion and willingness to learn and move forward step by step without backing out. 

\section{5. Portfolio Website Summary}

My portfolio website is hosted using GitHub Pages:

\href{https://adwaaaiith.github.io/Portfolio/}{https://adwaaaiith.github.io/Portfolio/}

The website contains the following sections:

\begin{itemize}[leftmargin=*]
    \item A home page with my name, photo, introduction, and CV link.
    \item An About section explaining what i'm and what i'm studying.
    \item An Education section with my university and school background.
    \item A Projects section which have my portfolio website and scientific calculator.
    \item A Skills section listing what all programming languages I am learning.
    \item A Contact section with my email, phone number, and location.
\end{itemize}

I changed the website structure from a normal down viewing page to a sidebar design. I done this to make it simple and look a bit different from others.

\section{6. GitHub Repository Structure}

The GitHub repository is available here:

\href{https://github.com/adwaaaiith/Portfolio}{https://github.com/adwaaaiith/Portfolio}

The repository contains:

\begin{itemize}[leftmargin=*]
    \item \texttt{index.html}: the main webpage portfolio.
    \item \texttt{style.css}: styling for the website.
    \item \texttt{script.js}: simple JavaScript for the calculator project page.
    \item \texttt{calculator.py}: the original Python scientific calculator.
    \item \texttt{calculator-project.html}: a simple page to show and try the calculator.
    \item \texttt{Adwaid-Sreeletha-Anil-CV.pdf}: the CV linked from the portfolio.
    \item \texttt{CV/}: a folder containing the LaTeX CV source files.
\end{itemize}

The \texttt{CV} folder includes the PDF, \texttt{resume.tex}, \texttt{resume.cls}, and \texttt{image.png}. This will make everyone to easily see my work and used codes.

\section{7. CV and Application Talking Points}

I created my CV using LaTeX. It includes my personal details, education, skills, projects, and links to my portfolio and GitHub repository. 

Important talking points from my application are:

\begin{itemize}[leftmargin=*]
    \item I am a beginner, but I am highly motivated and learning computer science from basics.
    \item I can create a simple responsive website using HTML and CSS.
    \item I can write and explain a beginner Python program.
    \item I can use GitHub to organize and publish project files.
    \item I can use LaTeX to create a professional CV.
\end{itemize}

\section{8. Python Project Overview}

The scientific calculator is a beginner Python project. It runs in the terminal and allows the user to select different mathematical operations. The calculator uses Python input, conditional statements, loops, and the \texttt{math} module.

Python code link:

\href{https://github.com/adwaaaiith/Portfolio/blob/main/calculator.py}{https://github.com/adwaaaiith/Portfolio/blob/main/calculator.py}

\subsection{Project Features}

\begin{itemize}[leftmargin=*]
    \item Addition, subtraction, multiplication, and division.
    \item Power and square root.
    \item Sine, cosine, and tangent.
    \item Logarithm and natural logarithm.
    \item Percentage conversion.
    \item Factorial calculation.
    \item Rounding numbers.
    \item Basic error handling for invalid inputs and division by zero.
\end{itemize}

\subsection{Technical Skills Demonstrated}

This project demonstrates beginner level Python knowledge, like:

\begin{itemize}[leftmargin=*]
    \item Using the \texttt{math} module.
    \item Taking user input.
    \item Using \texttt{if}, \texttt{elif}, and \texttt{else} statements.
    \item Using a loop to keep the program running.
    \item Checking for invalid input using \texttt{try} and \texttt{except}.
    \item Handling simple mathematical errors.
\end{itemize}

\section{9. Design and Implementation Decisions}

I chose a simple blue colour theme because blue is a simple colour.

The website uses a sidebar layout which makes it easier for the viewer to easily move between sections. I also kept the design simple for everyone to understand it easily.

The calculator project page is also kept simple. It shows my Python code and includes a small browser version so the viewer can understand how the calculator works without downloading.

\section{10. Strengths, Weaknesses, and Future Improvements}

\textbf{Strengths:}

\begin{itemize}[leftmargin=*]
    \item The portfolio is simple and easy to understand.
    \item The repository contains all the project files.
    \item The CV is available as a PDF and LaTeX format.
    \item The Python calculator shows basic programming logic clearly.
\end{itemize}

\textbf{Weaknesses:}

\begin{itemize}[leftmargin=*]
    \item The number of projects are very less.
    \item The calculator is a beginner terminal program, not a full application.
    \item The design is simple and could be improved later with more experience.
\end{itemize}

\textbf{Future Improvement:}

\begin{itemize}[leftmargin=*]
    \item Add more projects, such as a small web app.
    \item Improve the calculator with a graphical user interface.
    \item Continue improving JavaScript and Python skills.
\end{itemize}

\section{11. Conclusion}

This portfolio project helped me to understand how to make a personal portfolio, publish it on github, create CV in LaTeX and share my python programme.

\section{References}

\begin{itemize}
    \item GitHub Docs. (n.d.) \textit{GitHub Pages documentation}. Available at: \url{https://docs.github.com/en/pages}

    \item Python Software Foundation. (n.d.) \textit{Python documentation}. Available at: \url{https://docs.python.org/3/}

    \item Overleaf. (n.d.) \textit{LaTeX documentation and tutorials}. Available at: \url{https://www.overleaf.com/learn}
    \end{itemize}

\end{document}